\documentclass[11pt,a4paper]{article}
\usepackage[utf8]{inputenc}
\usepackage[T1]{fontenc}
\usepackage{lmodern}
\usepackage{textcomp}
\usepackage{amsmath,amssymb}
\usepackage{graphicx}
\usepackage{booktabs}
\usepackage{longtable}
\usepackage{array}
\usepackage{listings}
\usepackage{xcolor}
\usepackage[margin=1in]{geometry}
\usepackage{float}
\usepackage[hidelinks,breaklinks]{hyperref}
\usepackage{microtype}
\setlength{\parskip}{4pt}
\lstset{basicstyle=\ttfamily\scriptsize,breaklines=true,breakatwhitespace=false,
  postbreak=\mbox{\textcolor{gray}{$\hookrightarrow$}\space},frame=single,
  framerule=0.3pt,rulecolor=\color{gray!60},columns=fullflexible,keepspaces=true,
  showstringspaces=false,numbers=none,xleftmargin=4pt,xrightmargin=4pt,
  literate={\_}{{\_}}1}
\graphicspath{{figures/}}
\providecommand{\tightlist}{}
\title{Transcript: The 1.0--1.2 TeV Thermal Higgsino Against Planck, LZ, XENONnT and Gamma-Ray Limits: A Toy-Model Study}
\author{HEP-Genesis-Agent session record \\ \small Project: ramp-combo-higgsino-dd}
\date{2026-10-16}
\begin{document}
\maketitle
\noindent\textit{\small The conversation verbatim, generated mechanically from the session record.}
\vspace{1em}



\appendix
\section{Conversation transcript}
\label{app:transcript}
The exchange verbatim (assistant text as shown to the user; tool calls summarized by name --- see the tool-call log for arguments and results).
\subsection*{User \small(2026-10-16 14:31:17)}
Survey recent arXiv work on the higgsino thermal relic target (1.0-1.2 TeV). Use the inelastic-DM tools to compute the relic density at several masses; compare with LZ/XENONnT sensitivities and the Planck abundance. Where does the higgsino sit relative to current limits?

\subsection*{Assistant \small(research mode, 2026-10-16 14:31:17)}
\textit{Plan:}\begin{enumerate}
  \item (closed) Compute the thermal relic density Omega h\textasciicircum{}2 of a higgsino-like (inelastic) DM candidate at m = 0.9, 1.0, 1.1, 1.2 and 1.3 TeV with the inelastic-DM tools. Compare each value with Planck (Omega h\textasciicircum{}2 = 0.120 +/- 0.001) and locate the mass where the tool's relic curve crosses the Planck abundance.
  \item (closed) Place the thermal higgsino (m about 1.0-1.2 TeV) relative to LZ/XENONnT. Compute inelastic recoil spectra and expected event counts for delta = 300-600 keV, scan the (m\_chi, delta) plane against all constraints, and identify where the kinematic reach of xenon ends.
  \item (blocked) Verify and cite the key literature: the higgsino relic crossing mass, the LZ 2.84 t-yr extended-window result, the solar-capture delta bound, the elastic loop-level cross section relative to the XENONnT limits and the neutrino fog, and the Planck Omega h\textasciicircum{}2. Obtain verbatim quotes.
  \item (blocked) Cite the established, non-2026 baseline: the Planck 2018 Omega\_c h\textasciicircum{}2, the Sommerfeld-corrected pure-higgsino thermal mass from pre-2026 literature, and the XENONnT/LZ elastic SI limit at about 1 TeV compared with the loop-level higgsino cross section and the neutrino fog.
  \item (closed) Map, with the inelastic-DM tools, the mass splitting delta at which a thermal higgsino gives about 1 and about 0.3 expected LZ events. Test how sensitive that window is to the halo velocity cutoff, then compare it with the literature fits: 339 keV (2609.26698), about 350 keV (2609.01504) and 377 keV (2609.02608).
  \item (closed) Establish the current indirect-detection status of the 1.0-1.2 TeV thermal higgsino from the literature: the H.E.S.S., Fermi-LAT and MAGIC gamma-ray line and continuum limits against the Sommerfeld-enhanced s-wave annihilation cross section, and the CTAO projections. This replaces the server's p-wave treatment, which is physically wrong for a higgsino.
\end{enumerate}
\textit{Tool calls:} \texttt{list\_sterile\_nu\_skills}, \texttt{list\_observables}, \texttt{list\_agent\_files}, \texttt{list\_input\_files}, \texttt{list\_datasets}, \texttt{list\_files}, \texttt{get\_dispatch}, \texttt{search\_arxiv\_abstracts}, \texttt{search\_arxiv\_abstracts}, \texttt{search\_arxiv\_abstracts}, \texttt{search\_arxiv}, \texttt{web\_search}, \texttt{search\_arxiv}, \texttt{search\_arxiv}, \texttt{search\_arxiv}, \texttt{search\_arxiv}, \texttt{search\_arxiv}, \texttt{search\_arxiv}, \texttt{describe\_tools}, \texttt{describe\_tools}, \texttt{predict\_recoil}, \texttt{compute\_relic\_density}, \texttt{predict\_recoil}, \texttt{compute\_relic\_density}, \texttt{predict\_recoil}, \texttt{compute\_relic\_density}, \texttt{predict\_recoil}, \texttt{compute\_relic\_density}, \texttt{predict\_recoil}, \texttt{compute\_relic\_density}, \texttt{check\_all\_constraints}, \texttt{check\_all\_constraints}, \texttt{check\_all\_constraints}, \texttt{check\_all\_constraints}, \texttt{predict\_recoil}, \texttt{predict\_recoil}, \texttt{predict\_recoil}, \texttt{scan\_mass\_splitting}, \texttt{plot\_recoil\_spectrum}, \texttt{predict\_annual\_modulation}, \texttt{post\_finding}, \texttt{post\_finding}, \texttt{post\_question}, \texttt{post\_finding}, \texttt{post\_finding}, \texttt{post\_finding}, \texttt{plot\_constraint\_summary}, \texttt{post\_finding}, \texttt{get\_arxiv\_metadata}, \texttt{get\_arxiv\_metadata}, \texttt{download\_arxiv\_full\_text}, \texttt{download\_arxiv\_full\_text}, \texttt{download\_full\_arxiv\_paper}, \texttt{search\_arxiv}, \texttt{cite}, \texttt{search\_arxiv}, \texttt{search\_arxiv}, \texttt{post\_finding}, \texttt{post\_finding}, \texttt{post\_dead\_end}, \texttt{search\_arxiv}, \texttt{post\_question}, \texttt{predict\_recoil}, \texttt{predict\_recoil}, \texttt{predict\_recoil}, \texttt{predict\_recoil}, \texttt{predict\_recoil}, \texttt{predict\_recoil}, \texttt{search\_arxiv}, \texttt{search\_arxiv}, \texttt{get\_arxiv\_metadata}, \texttt{search\_arxiv}, \texttt{post\_finding}, \texttt{post\_finding}, \texttt{get\_arxiv\_metadata}, \texttt{download\_arxiv\_full\_text}, \texttt{post\_finding}, \texttt{post\_finding}, \texttt{get\_arxiv\_metadata}, \texttt{get\_arxiv\_metadata}, \texttt{download\_full\_arxiv\_paper}, \texttt{download\_arxiv\_full\_text}, \texttt{search\_arxiv}, \texttt{cite}, \texttt{search\_arxiv}, \texttt{search\_arxiv}, \texttt{post\_finding}, \texttt{post\_finding}, \texttt{post\_dead\_end}\par
\textbf{Survey} --- 52 tools on 5 servers, 7 read-only probes, 18 papers found.

Note first: I did not compute any relic densities. No inelastic-DM or relic-density tools were available to me, only literature search tools, so everything below comes from paper abstracts. Second note: the 'recent arXiv work' shown here was found through the OpenAlex-backed search\_arxiv tool, which returned papers dated June to October 2026. These IDs should be checked on arxiv.org before anyone cites them. The direct arXiv API search\_arxiv\_abstracts tool returned HTTP 503/429 errors. On the physics: when the Sommerfeld effect and coannihilation are included, a pure higgsino gives the full Planck abundance at m about 1.0-1.2 TeV. The cited values are about 1.08-1.1 TeV (2609.02775, 2609.02608), about 1.2 TeV (2410.15733) and an upper end near 1150 GeV (2403.19598). Lighter thermal higgsinos are under-abundant, typically by a factor of 5-20 in natural SUSY (1803.11210). For direct detection, the pure higgsino is inelastic dark matter: its Z-mediated scattering requires a neutral-state mass splitting delta of up to about a few hundred keV. If delta is larger than the kinematic reach of xenon, elastic loop-level scattering sits near or below the neutrino fog (2609.01504, 2410.02723). Higgsino-gaugino mixing raises the elastic rate and is what lets direct detection probe the thermal target (2606.31013). In 2026, LZ released a 2.84 tonne-year search with an extended nuclear-recoil window (2609.02823). It found one event at 248 keV, with 2.6 sigma global and 3.4 sigma local significance. A roughly 1.1 TeV higgsino with delta of about 340-380 keV under the Standard Halo Model, or about 480-520 keV if the Large Magellanic Cloud's effect on the halo velocity tail is included, can produce this event. However, IceCube solar-capture neutrino limits require delta \textgreater{} 566 keV (2609.02775, confirmed by 2609.11833 and 2609.07807; 2609.06760 quotes 0.51-0.56 MeV). Two further problems: the empty sideband above 350 keV is in tension with the thermal-mass fit (2609.09385, 2609.26698), and halo uncertainties are not enough to close the gap (2609.15321). Bottom line: the full-abundance thermal higgsino at about 1.1 TeV is not excluded as a model. Its sub-MeV-splitting interpretation of the LZ event is disfavoured by solar capture. Splittings of order an MeV or more remain allowed (2610.00226) and leave it effectively out of reach of current xenon searches. Indirect gamma-ray line searches (CTA/CTAO, 2008.00692) are the main future test.

Quantities worth checking: Thermal relic Omega h\textasciicircum{}2 versus higgsino mass at about 0.9, 1.0, 1.1, 1.2 and 1.3 TeV, with and without the Sommerfeld effect. Compare with the Planck value Omega h\textasciicircum{}2 of about 0.120 +/- 0.001 (standard Planck 2018 value; the Planck paper was not retrieved here). The literature crossing point is about 1.08-1.2 TeV.; Kinematic maximum splitting delta\_max = mu v\textasciicircum{}2/2 for xenon (and argon) at m = 1.1 TeV, for v\_esc of about 544 km/s in the Earth frame (about 544+250 km/s).; Endothermic recoil spectrum dR/dE\_R for m = 1.0-1.2 TeV and delta = 300-600 keV, with Helm form factor. Count events in [14,225), [225,271] and (271,800] keV for 2.84 tonne-years.; Z-mediated inelastic cross section per nucleon for a pure higgsino (vector coupling, roughly 1e-39 cm\textasciicircum{}2 scale) versus the rate needed to give 1 event.; Elastic loop-level higgsino-nucleon SI cross section (about 1e-48 to 1e-49 cm\textasciicircum{}2) versus the LZ/XENONnT elastic limits at 1.1 TeV and the neutrino fog.; Solar-capture threshold delta of about 566 keV versus the LZ-preferred delta (about 340-380 keV for the Standard Halo Model, about 480-520 keV with the LMC)..

\textbf{Plan} --- 4 lines of investigation:

\begin{enumerate}
  \item Compute the thermal relic density Omega h\textasciicircum{}2 of a higgsino-like (inelastic) DM candidate at m = 0.9, 1.0, 1.1, 1.2 and 1.3 TeV with the inelastic-DM tools. Compare each value with Planck (Omega h\textasciicircum{}2 = 0.120 +/- 0.001) and locate the mass where the tool's relic curve crosses the Planck abundance. --- This answers the core quantitative part of the task directly. The survey computed no relic densities. Its literature crossing point of about 1.08-1.2 TeV needs a tool-based cross-check, and we need to know which physics the tool includes (Sommerfeld effect, coannihilation, higgsino couplings).
  \item Place the thermal higgsino (m about 1.0-1.2 TeV) relative to LZ/XENONnT. Compute inelastic recoil spectra and expected event counts for delta = 300-600 keV, scan the (m\_chi, delta) plane against all constraints, and identify where the kinematic reach of xenon ends. --- This answers 'where does the higgsino sit relative to current limits'. The survey's key claims need tool-based checks: the LZ 248 keV event is fit by delta of about 340-380 keV, solar capture requires delta \textgreater{} 566 keV, and large-delta higgsinos fall out of xenon reach.
  \item Verify and cite the key literature: the higgsino relic crossing mass, the LZ 2.84 t-yr extended-window result, the solar-capture delta bound, the elastic loop-level cross section relative to the XENONnT limits and the neutrino fog, and the Planck Omega h\textasciicircum{}2. Obtain verbatim quotes. --- The survey's 2026-dated IDs came from an OpenAlex-backed search. They must be confirmed on arXiv before anything is cited, and the Planck value was never retrieved. The tool results from d-01 and d-02 need a validated literature baseline to compare against.
  \item Cite the established, non-2026 baseline: the Planck 2018 Omega\_c h\textasciicircum{}2, the Sommerfeld-corrected pure-higgsino thermal mass from pre-2026 literature, and the XENONnT/LZ elastic SI limit at about 1 TeV compared with the loop-level higgsino cross section and the neutrino fog. --- The original d-03 cannot cover these within 13 calls. These anchors are independent of whether the 2026 IDs resolve, so they protect the final answer if the 2026 papers prove unverifiable. (added by the critic)
Critic's view: revise --- Over budget. As written it needs about 15 calls: 1 describe + 5 relic + 3 Sommerfeld toggles + 3 delta checks + 2 check\_all. Drop the delta-sensitivity step. Keep the Sommerfeld toggle only if describe\_tools shows a switch exists. This brings it to about 9-11 calls.; The serve-idm server is probably a generic inelastic-DM model, not a pure higgsino with SU(2) couplings, Sommerfeld enhancement and full coannihilation. If describe\_tools does not state higgsino electroweak couplings and Sommerfeld enhancement, label the results as a toy model. In that case, do not report a 'higgsino Planck-crossing mass'. Report only the tool's crossing mass and how far it sits from the 1.0-1.2 TeV literature target. Quote the tool's model description verbatim.; The delta-sensitivity step (0.3, 0.566, 1.0 MeV) is physically uninformative. A sub-MeV splitting is far below T\_f of about m/20, which is about 50 GeV, so the relic density should not depend on delta. If the tool does show delta dependence, that indicates a modelling artifact and should be reported as one. Check whether delta is entered in keV or MeV before calling..

\end{enumerate}
\textbf{Wave 1} --- now exploring, in parallel: Compute the thermal relic density Omega h\textasciicircum{}2 of a higgsino-like (inelastic) DM candidate at m = 0.9, 1.0, 1.1, 1.2 and 1.3 TeV with the inelastic-DM tools. Compare each value with Planck (Omega h\textasciicircum{}2 = 0.120 +/- 0.001) and locate the mass where the tool's relic curve crosses the Planck abundance.; Place the thermal higgsino (m about 1.0-1.2 TeV) relative to LZ/XENONnT. Compute inelastic recoil spectra and expected event counts for delta = 300-600 keV, scan the (m\_chi, delta) plane against all constraints, and identify where the kinematic reach of xenon ends.; Cite the established, non-2026 baseline: the Planck 2018 Omega\_c h\textasciicircum{}2, the Sommerfeld-corrected pure-higgsino thermal mass from pre-2026 literature, and the XENONnT/LZ elastic SI limit at about 1 TeV compared with the loop-level higgsino cross section and the neutrino fog..

\textbf{Finding:} The inelastic-DM server's compute\_relic\_density (is\_higgsino=true) gives Omega h\textasciicircum{}2 = 0.0731, 0.0898, 0.1083, 0.1285, 0.1503 at m\_chi = 900, 1000, 1100, 1200, 1300 GeV. Linear interpolation crosses Planck 0.120 at m about 1158 GeV (+/-6 GeV for the +/-0.0012 band). A direct check at 1159 GeV returns Omega h\textasciicircum{}2 = 0.120. This lies within the 1.08-1.2 TeV literature range but about 60-80 GeV above the 1.08-1.1 TeV values.

\textbf{Finding:} check\_all\_constraints (is\_higgsino) at m=1100 GeV with delta = 400, 566 and 1000 keV, and at m=1159 GeV with delta=400 keV, returns VIABLE every time. Fermi-LAT and CMB are not excluded, sigma\_n = 7.43e-39 cm\textasciicircum{}2, and N\_expected\_LZ = 0.00 at every point. N\_LZ = 0 at delta = 400 keV contradicts the literature, which reports about 1 event at delta about 340-380 keV. The server's present-day sigma\_v (about 7e-42 to 5e-41 cm\textasciicircum{}3/s, labelled 'p-wave + mass-splitting suppressed') is also far below the s-wave higgsino value of about 1e-26 cm\textasciicircum{}3/s.

\textbf{Completed --- Compute the thermal relic density Omega h\textasciicircum{}2 of a higgsino-like (inelastic) DM candidate at m = 0.9, 1.0, 1.1, 1.2 and 1.3 TeV with the inelastic-DM tools. Compare each value with Planck (Omega h\textasciicircum{}2 = 0.120 +/- 0.001) and locate the mass where the tool's relic curve crosses the Planck abundance..} Relic densities for a higgsino-like dark matter candidate (compute\_relic\_density with is\_higgsino=true; the tool sets the annihilation cross section automatically): Omega h\textasciicircum{}2 = 0.0731, 0.0898, 0.1083, 0.1285 and 0.1503 at m = 900, 1000, 1100, 1200 and 1300 GeV (tc-...-23/25/27/29/31). The curve reaches the Planck value of 0.120 at about 1158 GeV, +/-6 GeV from the server's +/-0.0012 Planck band (the brief's +/-0.001 gives about +/-5 GeV). A direct call at 1159 GeV returns 0.120 (tc-...-33). This is inside the 1.08-1.2 TeV literature range. It is about 60-80 GeV above the 1.08-1.1 TeV values (2609.02775, 2609.02608) and close to the about 1.2 TeV value (2410.15733), so agreement is only partial.

Physics in the tool: the higgsino setting is a simple scaling of the cross section, sigma\_v = 2.5e-26 cm\textasciicircum{}3/s x (1100 GeV/m)\textasciicircum{}2, fed into a Kolb-Turner freeze-out estimate (x\_f about 24.5, g* = 86.25). There is no Sommerfeld or coannihilation switch; the only inputs are m\_chi\_GeV, is\_higgsino and sigma\_v\_cm3\_s. As a result, the with/without-Sommerfeld comparison (step 3) and the splitting-sensitivity runs (step 4) could not be done with compute\_relic\_density. Its closeness to the literature reflects how it was calibrated, not a full calculation.

Constraint checks (check\_all\_constraints): at 1100 GeV with delta = 400, 566 and 1000 keV, and at 1159 GeV with delta = 400 keV, every point is rated VIABLE. In each case Fermi-LAT and Planck CMB show no exclusion, sigma\_n = 7.43e-39 cm\textasciicircum{}2 and N\_expected\_LZ = 0.00 (tc-...-32/33/34/35). The relic verdict passes at 1100 GeV despite the tool's own 9.8 sigma tension.

Caveats:

\begin{itemize}
  \item The pass/consistent flags are True even at 7-25 sigma tension, which looks like a bug.
  \item N\_LZ = 0 at delta = 400 keV conflicts with the literature, where splittings of about 340-380 keV give about one event (posted as q-06).
  \item The indirect-detection treatment assumes p-wave annihilation suppressed by the mass splitting, giving sigma\_v of about 1e-41 cm\textasciicircum{}3/s today. That is inappropriate for an s-wave, Sommerfeld-enhanced higgsino.
\end{itemize}
Posted as findings f-001 and f-002 and question q-06. No figures or files were produced.

\textbf{Finding:} With the inelastic-DM server (Z-exchange higgsino, sigma\_n = 7.43e-39 cm\textasciicircum{}2, LZ 2.84 t-yr extended NR window, v\_esc+v\_E = 776 km/s), the expected LZ count at m\_chi = 1100 GeV is 744.6 at delta = 300 keV, 1.97 at 350 keV (P(N\textgreater{}=1) = 0.86, peak E\_R = 315 keV), and 0.00 at 380, 500, 566 and 700 keV (peak E\_R = 342, 450, 510, 630 keV). The 248 keV event needs v\_min = 762 km/s at 350 keV (accessible) and 799 km/s at 380 keV (not accessible).

\textbf{Finding:} Xenon's kinematic reach at the thermal higgsino mass is delta\_max = mu v\_max\textasciicircum{}2/2. With m\_Xe of about 122 GeV, m\_chi = 1100 GeV gives mu = 109.8 GeV. For v\_max = 776 km/s (the tool's value, beta\textasciicircum{}2 = 6.70e-6) this gives delta\_max of about 368 keV; for v\_max = 794 km/s (544+250) it gives about 385 keV. This agrees with the tool: 350 keV is accessible and 380 keV gives 0 events. At delta = 380 keV the 248 keV v\_min is 802, 799 and 796 km/s for m = 1000, 1100 and 1200 GeV, all inaccessible, so the reach depends only weakly on mass.

\textbf{Finding:} The scan over m\_chi = 800-1400 GeV and delta = 0-800 keV (17x17 grid) finds 144 of 289 points passing all constraints and 9 points that predict 0.3-5 LZ events. The 1100 GeV, 380 keV recoil-spectrum plot is at /Users/nramachandra/github-integration/inelastic-dm-server/output/recoil\_spectrum.png. The annual modulation at 1100 GeV, 350 keV and E\_R = 248 keV has a fractional amplitude of 0.985, peaking on June 2; the event date of September 1 sees 22.5\% of the peak rate.

\textbf{Finding:} plot\_constraint\_summary on the 17x17 scan (m\_chi = 800-1400 GeV, delta = 0-800 keV) wrote /Users/nramachandra/github-integration/inelastic-dm-server/output/constraint\_summary.png and reports the best-fit point at m\_chi = 988 GeV, delta = 350 keV.

\textbf{Completed --- Place the thermal higgsino (m about 1.0-1.2 TeV) relative to LZ/XENONnT. Compute inelastic recoil spectra and expected event counts for delta = 300-600 keV, scan the (m\_chi, delta) plane against all constraints, and identify where the kinematic reach of xenon ends..} (Open items: (1) The tool gives only total counts, with no breakdown into the [14,225), [225,271] and (271,800] keV bins. (2) The scan and plot files went to the server's default directory, not the run directory, because the tools have no output\_dir parameter. (3) The 566 keV and literature 340-380 keV bands were not drawn on the plot; the plotting tool has no overlay option. (4) Halo-parameter sensitivity (q-04) was not tested, because the tool exposes no v0 or v\_esc parameters.) The inelastic-DM server models the higgsino with Z exchange (sigma\_n = 7.43e-39 cm\textasciicircum{}2), the LZ 2.84 t-yr extended nuclear-recoil window, and a halo cutoff of v\_esc + v\_E = 776 km/s (tc-...-20, -21). At m\_chi = 1100 GeV, the expected LZ event counts are: 744.6 at delta = 300 keV (grossly excluded against 1 observed event); 1.97 at 350 keV (P(N\textgreater{}=1) = 0.86, peak recoil energy E\_R = 315 keV); and 0.00 at 380, 500, 566 and 700 keV, where the spectrum would peak at 342, 450, 510 and 630 keV (finding f-003). At delta = 380 keV the result does not depend much on mass: the minimum speed needed for the 248 keV event is v\_min = 802, 799 and 796 km/s at m = 1000, 1100 and 1200 GeV, all above 776 km/s and so out of reach. Xenon's kinematic reach, worked by hand: delta\_max = mu v\_max\textasciicircum{}2/2, with mu = 109.8 GeV. That gives about 368 keV for v\_max = 776 km/s and about 385 keV for 544+250 km/s, consistent with the tool's cutoff between 350 and 380 keV (f-004). The scan over m = 800-1400 GeV and delta = 0-800 keV (17x17 points) found 144 of 289 points viable and 9 points predicting 0.3-5 LZ events. Its best fit is m = 988 GeV, delta = 350 keV. Files: /Users/nramachandra/github-integration/inelastic-dm-server/output/scan\_idm.npz and constraint\_summary.png (f-005, f-006). The recoil spectrum plot for 1100 GeV, 380 keV is at /Users/nramachandra/github-integration/inelastic-dm-server/output/recoil\_spectrum.png. Annual modulation at 1100 GeV, 350 keV and E\_R = 248 keV: 98.5\% amplitude, peaking on June 2; the event date of September 1 sees 22.5\% of the peak rate. Regions at 1.1 TeV: delta below about 340 keV is excluded by xenon (hundreds of events predicted); delta of about 345-368 keV is allowed and can explain the 248 keV event (tool value); delta above about 370 keV is out of xenon's reach. The literature solar-capture bound (delta \textgreater{} 566 keV, 2609.02775) therefore puts every solar-capture-allowed thermal higgsino out of LZ inelastic reach under the Standard Halo Model.

\textbf{Finding:} The pre-2026 Sommerfeld- and coannihilation-corrected thermal pure-higgsino mass is about 1.2 TeV. The abstract of arXiv:2410.15733 (Fukuda, Li, Moroi, Niki, JHEP06(2025)091, DOI 10.1007/jhep06(2025)091) says: ''in which the observed DM abundance is realized with the Higgsino mass of \textasciitilde{}1.2 TeV''.

\textbf{Finding:} XENONnT's 3.1 tonne-year elastic WIMP search (arXiv:2502.18005, DOI 10.1103/msw4-t342) found no significant excess. Its spin-independent limit has a minimum of 1.7e-47 cm\textasciicircum{}2 (90\% CL) at 30 GeV, with best median sensitivity 1.4e-47 cm\textasciicircum{}2 at 41 GeV. The earlier 1.09 t-yr result (arXiv:2303.14729) reached 2.58e-47 cm\textasciicircum{}2 at 28 GeV. The abstracts do not give the limit at 1 TeV. Because the limit scales roughly linearly with mass at high mass, it should be about 1e-45 cm\textasciicircum{}2 near 1 TeV, but that is an inference, not a quoted value.

\textbf{Blocked --- Cite the established, non-2026 baseline: the Planck 2018 Omega\_c h\textasciicircum{}2, the Sommerfeld-corrected pure-higgsino thermal mass from pre-2026 literature, and the XENONnT/LZ elastic SI limit at about 1 TeV compared with the loop-level higgsino cross section and the neutrino fog..} (The arXiv-API-backed tools (get\_arxiv\_metadata, download\_arxiv\_full\_text, download\_full\_arxiv\_paper and cite's metadata check) returned 'not found', 429 or timeout errors even for 1807.06209 and 2410.15733 (tc-...-19, -50, -51, -52, -53, -55). So cite() cannot validate any quote and no full text can be read. Only OpenAlex abstracts were available.) Two of the three anchors were confirmed from abstracts returned by the OpenAlex-backed search\_arxiv tool. None could be added to references.bib with cite(), so the success criterion of 3 cite() entries was not met. (1) Thermal higgsino mass, finding f-007: arXiv:2410.15733 (Fukuda, Li, Moroi, Niki, JHEP06(2025)091, DOI 10.1007/jhep06(2025)091) states 'in which the observed DM abundance is realized with the Higgsino mass of \textasciitilde{}1.2 TeV', with the Sommerfeld effect and coannihilation included (tc-261016\_1-w1-d-04-54). (2) XENONnT elastic limits, finding f-008: arXiv:2502.18005 (3.1 t-yr exposure, DOI 10.1103/msw4-t342) gives a minimum spin-independent limit of 1.7e-47 cm\textasciicircum{}2 at 30 GeV (90\% CL) and a best median sensitivity of 1.4e-47 cm\textasciicircum{}2 at 41 GeV. arXiv:2303.14729 (1.09 t-yr) gives 2.58e-47 cm\textasciicircum{}2 at 28 GeV (tc-261016\_1-w1-d-04-57). The limit at 1 TeV is not quoted in either abstract. The figure of about 1e-45 cm\textasciicircum{}2 is only an extrapolation from linear high-mass scaling. (3) Planck Omega\_c h\textasciicircum{}2 = 0.120 +/- 0.001 (1807.06209) was NOT retrieved: get\_arxiv\_metadata and both download tools returned 'not found' or failed, and a targeted search did not surface the Planck VI paper. The loop-level higgsino cross section (about 1e-48 to 1e-49 cm\textasciicircum{}2) and the neutrino-fog comparison were not sourced. Dead end x-01 and question q-07 are posted.

\textbf{Critic, wave 1} --- continue: 5 findings stand up to the evidence, 3 questioned ('check\_all\_constraints (is\_higgsino) at m=1100 GeV with delta = 400, 566 and 1000 keV, and ': The excerpts for tc-32/33/34/35 are cut off before the overall verdict and the CMB block, so 'VIABLE every time' and 'CMB not excluded' are not shown. The excerpts also show consistent\_with\_1\_event = False next to N\_LZ = 0, which works against a clean VIABLE reading. The claim that N\_LZ = 0 at delta = 400 keV contradicts the literature is overstated: 400 keV lies above the quoted 340-380 keV band. The real tension is at 380 keV (tool gives 0 events, 2609.02608 fits about 377 keV), and that point belongs to f-003 and f-004, not here. The remaining numbers (sigma\_n, present-day sigma\_v of 6.7e-42 to 5.2e-41, Fermi not excluded) do match.; 'The pre-2026 Sommerfeld- and coannihilation-corrected thermal pure-higgsino mass is about ': The quoted excerpt of tc-261016\_1-w1-d-04-54 shows only Paper 1, arXiv:1404.0682 (Low and Wang). It does not contain 2410.15733 or the quoted sentence 'Higgsino mass of \textasciitilde{}1.2 TeV'. The evidence as posted does not contain the claimed value.; 'XENONnT's 3.1 tonne-year elastic WIMP search (arXiv:2502.18005, DOI 10.1103/msw4-t342) fou': The excerpt of tc-261016\_1-w1-d-04-57 shows only the title and author list of 2303.14729. None of the claimed limits (1.7e-47 at 30 GeV, 1.4e-47 at 41 GeV, 2.58e-47 at 28 GeV) appear, and 2502.18005 is not visible. The claimed values are not in the evidence.).

Issues: The arithmetic checks out: sigma\_v = 2.5e-26 x (1100/m)\textasciicircum{}2 reproduces 3.73e-26 at 900 GeV, and interpolating between 1100 and 1200 GeV gives about 1158 GeV. However, the tool is a single-channel Kolb-Turner scaling calibrated to 1.1 TeV. The 1158 GeV crossing is a calibration artefact, not a higgsino prediction. The final report must label it as a toy-model crossing, as critique c-01 required.; The scan counts 144 of 289 points as viable. That count uses the same relic 'passes' flag that returns True at 7-25 sigma tension, so it probably includes points that are far from the Planck abundance. It also relies on the p-wave-suppressed indirect-detection model, which is wrong for a higgsino. The viable fraction should not be quoted as physics until it is recomputed with a correct relic criterion.; The best fit at m = 988 GeV, delta = 350 keV fits the LZ event, not the thermal relic. With the tool's own relic curve, 988 GeV gives Omega h\textasciicircum{}2 of about 0.087, which is under-abundant. The best fit is therefore not a full-abundance thermal higgsino, and this should be stated..

\textbf{Next steps after wave 1} --- One wave remains, with 19 calls, so the plan is three directions sized at about 6, 7 and 5 calls. d-03, which is open and has not run, carries over first. It already covers the critic's literature-repair suggestion: re-capturing the abstracts behind f-007 and f-008 so the excerpts show the quoted numbers, sourcing the elastic loop-level cross section and the neutrino fog at 1 TeV, retrying the Planck value, and resolving the 2026 IDs (q-03). For that reason I add no duplicate literature direction. d-03 should cap itself at about 6 calls and drop cite() attempts if the arXiv API fails again. The first new direction addresses the critic's highest-severity issue: the delta window for one LZ event and its dependence on the velocity cutoff. That window sets where the thermal higgsino sits relative to xenon reach, between the kinematic edge near 370-385 keV and the solar-capture bound at 566 keV. The second new direction closes the unaddressed indirect-detection leg with sourced H.E.S.S./Fermi/CTAO status, replacing the tool's invalid p-wave treatment. The critic's suggested full rescan with a corrected relic criterion is not run, because it would not fit the budget. Instead, the final report should note that along the thermal line only delta is free: f-005's 144/289 viable count must not be quoted, and f-006's best fit at 988 GeV is under-abundant, so it is not a thermal higgsino. The 1158 GeV crossing must be labelled a calibrated toy-model value. The task does not ask for a facility (HPC) run, so none is scheduled.

New lines of investigation: Map, with the inelastic-DM tools, the mass splitting delta at which a thermal higgsino gives about 1 and about 0.3 expected LZ events. Test how sensitive that window is to the halo velocity cutoff, then compare it with the literature fits: 339 keV (2609.26698), about 350 keV (2609.01504) and 377 keV (2609.02608).; Establish the current indirect-detection status of the 1.0-1.2 TeV thermal higgsino from the literature: the H.E.S.S., Fermi-LAT and MAGIC gamma-ray line and continuum limits against the Sommerfeld-enhanced s-wave annihilation cross section, and the CTAO projections. This replaces the server's p-wave treatment, which is physically wrong for a higgsino..

\textbf{Wave 2} --- now exploring, in parallel: Verify and cite the key literature: the higgsino relic crossing mass, the LZ 2.84 t-yr extended-window result, the solar-capture delta bound, the elastic loop-level cross section relative to the XENONnT limits and the neutrino fog, and the Planck Omega h\textasciicircum{}2. Obtain verbatim quotes.; Map, with the inelastic-DM tools, the mass splitting delta at which a thermal higgsino gives about 1 and about 0.3 expected LZ events. Test how sensitive that window is to the halo velocity cutoff, then compare it with the literature fits: 339 keV (2609.26698), about 350 keV (2609.01504) and 377 keV (2609.02608).; Establish the current indirect-detection status of the 1.0-1.2 TeV thermal higgsino from the literature: the H.E.S.S., Fermi-LAT and MAGIC gamma-ray line and continuum limits against the Sommerfeld-enhanced s-wave annihilation cross section, and the CTAO projections. This replaces the server's p-wave treatment, which is physically wrong for a higgsino..

\textbf{Finding:} With predict\_recoil (Z-exchange higgsino, sigma\_n = 7.43e-39 cm\textasciicircum{}2, LZ 2.84 t-yr extended NR window, hard cutoff v\_esc+v\_E = 776 km/s), m\_chi = 1100 GeV gives N\_LZ = 1.97, 0.494, 0.085, 0.009, 0.00, 0.00 at delta = 350, 355, 360, 365, 370, 375 keV (peak E\_R = 315, 320, 324, 329, 333, 338 keV). Log-linear interpolation gives delta(N=1) = 352.5 keV and delta(N=0.3) = 356.4 keV, so the window is only about 4 keV wide (N falls about 4-6x per 5 keV). The 248 keV event stops being accessible between 360 and 365 keV (v\_min = 774 and 781 km/s). At m\_chi = 1159 GeV (the tool's thermal mass), delta = 360 keV gives N = 0.135 against 0.085 at 1100 GeV. That moves delta(N=1) up by only about 1-2 keV, to about 354 keV.

\textbf{Finding:} This answers q-06. predict\_recoil has no halo parameters (inputs: m\_chi\_GeV, delta\_keV, sigma\_n\_cm2, is\_higgsino), so v\_esc cannot be varied. The N\_LZ = 0 at delta \textgreater{}= 370 keV comes from the tool's hard 776 km/s cutoff, which puts delta\_max = mu v\_max\textasciicircum{}2/2 at about 368 keV. An analytic rescaling by (794/776)\textasciicircum{}2 = 1.047 moves the edge to about 385 keV, and moves the tool's delta(N=1) and delta(N=0.3) from 352.5 and 356.4 keV to about 369 and 373 keV. Against the literature: 350 keV (2609.01504) matches the tool directly (N = 1.97). The 776 vs 794 km/s cutoff accounts for about 17 keV of the 24.5 keV offset from 377 keV (2609.02608). The remaining roughly 5-8 keV has not been explained; plausible causes, none tested, are a smooth truncated-Maxwellian tail, a different v\_E at the event date, and differences in cross-section and efficiency normalization. The 339 keV value at 1 TeV (2609.26698) is a normalization-free sideband-ratio optimum. The tool would predict far more than 1 event there (744.6 at 300 keV), so 339 keV is not a direct comparison for the absolute count.

\textbf{Completed --- Map, with the inelastic-DM tools, the mass splitting delta at which a thermal higgsino gives about 1 and about 0.3 expected LZ events. Test how sensitive that window is to the halo velocity cutoff, then compare it with the literature fits: 339 keV (2609.26698), about 350 keV (2609.01504) and 377 keV (2609.02608)..} I mapped the window with predict\_recoil (Z-exchange higgsino, sigma\_n = 7.43e-39 cm\textasciicircum{}2, LZ 2.84 t-yr extended NR window, hard cutoff v\_esc+v\_E = 776 km/s). At m\_chi = 1100 GeV the expected LZ count is N\_LZ = 1.97, 0.494, 0.085, 0.009, 0.00 and 0.00 at delta = 350, 355, 360, 365, 370 and 375 keV, with the spectrum peaking at E\_R = 315-338 keV (tc-...-65 to -69). Interpolating in log N gives delta(N=1) = 352.5 keV and delta(N=0.3) = 356.4 keV at 1100 GeV, so the window is only about 4 keV wide. The 248 keV event becomes inaccessible between 360 and 365 keV, where its v\_min rises from 774 to 781 km/s. Mass dependence along the thermal line is weak: at 1159 GeV, delta = 360 keV gives N = 0.135 against 0.085 at 1100 GeV (tc-...-70), which moves delta(N=1) up by about 1-2 keV. The halo cutoff could not be tested with the tool, because predict\_recoil takes no halo parameters. Scaling delta\_max with v\_max\textasciicircum{}2 by hand, (794/776)\textasciicircum{}2 = 1.047, moves the kinematic edge from about 368 to about 385 keV and puts delta(N=1) near 369 keV and delta(N=0.3) near 373 keV. This answers q-06: N = 0 at delta \textgreater{}= 370 keV is caused by the hard 776 km/s cutoff, not by the energy window. Against the literature: 350 keV (2609.01504) agrees with the tool directly. The 776 vs 794 km/s cutoff accounts for about 17 keV of the 24.5 keV offset from 377 keV (2609.02608). The remaining roughly 5-8 keV is unexplained; possible causes such as the shape of the halo tail, the value of v\_E and the normalization were not tested. The 339 keV value at 1 TeV (2609.26698) comes from a sideband-ratio optimum that does not depend on the overall rate. With the tool's normalization, splittings at or below 340 keV give hundreds of events (744.6 at 300 keV), so 339 keV cannot be compared directly with an absolute event count. Findings f-009 and f-010 are posted. No figures or files were produced. I used 6 compute calls.

\textbf{Finding:} The indirect-detection status of the about 1.1 TeV thermal higgsino, taken from abstracts only: it is not excluded by current gamma-ray searches, and the result depends on the halo profile. (a) arXiv:2405.13104 (Rodd, Safdi, Xu, PRD 110, 043003, 2024) ran a 15-yr Fermi-LAT line-like search that comes 'within an order one factor' of the sensitivity needed to detect the higgsino, for some Milky Way analog profiles. In that analysis H.E.S.S. is subleading to Fermi for the higgsino. (b) arXiv:2210.03140 (Montanari, Moulin, Rodd, PRD 107, 043028) finds the higgsino 'out of reach by at least a factor of a few' even at the ultimate sensitivity of current H.E.S.S.-like IACTs. (c) arXiv:2608.07234 (H.E.S.S. Inner Galaxy Survey, 546 h, DOI 10.1103/d8tj-55kc, dated 2026 and unverified) quotes a line limit of 2.3e-28 cm\textasciicircum{}3/s at 1 TeV for an Einasto profile and states that thermal higgsino DM is 'probed for DM Milky Way models' for the first time. The abstract does not say whether it excludes the higgsino, or for which profiles.

\textbf{Finding:} CTAO is the decisive future test. arXiv:2008.00692 (Rinchiuso et al., PRD 103, 023011) finds that CTA 'has the potential to become the first experiment that has sensitivity to the thermal' higgsino. arXiv:2405.13104 finds that CTA and SWGO could detect the line-like signal of the 1.1 TeV higgsino 'at high significance... for all but the most pessimistic DM profiles'. arXiv:2506.08084 (PoS ICRC2025) projects that CTAO-North, observing at large zenith angles, could reach the required sensitivity within the next decade. The server's indirect-detection treatment (p-wave, mass-splitting suppressed, sigma\_v of about 1e-41 cm\textasciicircum{}3/s, f-002) is SUPERSEDED. A higgsino annihilates through an s-wave, Sommerfeld-enhanced channel, and the literature treats the thermal target as near the current Fermi/IACT reach rather than 15 orders of magnitude below it.

\textbf{Completed --- Establish the current indirect-detection status of the 1.0-1.2 TeV thermal higgsino from the literature: the H.E.S.S., Fermi-LAT and MAGIC gamma-ray line and continuum limits against the Sommerfeld-enhanced s-wave annihilation cross section, and the CTAO projections. This replaces the server's p-wave treatment, which is physically wrong for a higgsino..} Verdict, based on abstracts only (download\_arxiv\_full\_text on 2405.13104 returned 'not found', tc-261016\_1-w2-d-06-76): the 1.0-1.2 TeV thermal higgsino is not excluded by current gamma-ray line or continuum searches, and the result depends on the halo profile. Current searches: arXiv:2405.13104 (Rodd, Safdi, Xu, PRD 110, 043003) ran a 15-yr Fermi-LAT line-like search that comes 'within an order one factor' of the sensitivity needed to detect the higgsino, for some Milky Way analog density profiles from hydrodynamic simulations. The same paper finds H.E.S.S. subleading to Fermi for the higgsino. arXiv:2210.03140 (Montanari, Moulin, Rodd, PRD 107, 043028) finds the higgsino 'out of reach by at least a factor of a few' even at the ultimate sensitivity of current H.E.S.S.-like IACTs; the dominant systematic is the dark matter distribution in the inner Galaxy. The 2026-dated H.E.S.S. Inner Galaxy Survey, arXiv:2608.07234 (546 h, DOI 10.1103/d8tj-55kc, still to be checked on arxiv.org), quotes a line limit of 2.3e-28 cm\textasciicircum{}3/s at 1 TeV for an Einasto profile and says thermal higgsino DM is 'probed for DM Milky Way models' for the first time. Its abstract does not say whether it excludes the higgsino, or for which profiles. No MAGIC-specific higgsino limit came up in the searches. Future: arXiv:2008.00692 says CTA could be the first experiment sensitive to the thermal higgsino; 2405.13104 says CTA and SWGO could detect it 'at high significance' for all but the most pessimistic profiles; 2506.08084 projects CTAO-North could reach the required sensitivity within a decade. The server's indirect-detection treatment (p-wave, mass-splitting suppressed, sigma\_v of about 1e-41 cm\textasciicircum{}3/s, f-002) is explicitly flagged as superseded, because higgsino annihilation is s-wave and Sommerfeld-enhanced. The 'no Fermi-LAT exclusion' flags in the f-005 viable count are therefore not physically meaningful for indirect detection. The predicted line cross section at 1.1 TeV, and its ratio to the limit, could not be extracted without a full text. Findings f-011 and f-012 are posted (evidence tc-261016\_1-w2-d-06-63, -71, -72). No files or figures were produced.

\textbf{Finding:} The OpenAlex record for arXiv:2609.02823 (LZ, 'Search for dark matter particle interactions in an extended nuclear recoil energy window with the LUX-ZEPLIN (LZ) experiment', 2026-09-02, DOI 10.48550/arxiv.2609.02823) carries this abstract text: ''an exposure of 2.84 tonne-years'', ''An extended nuclear recoil energy window up to approximately 270 keV'', ''one event with characteristics consistent with a nuclear recoil of 248$\pm$23 (stat)$\pm$23 (sys) keV'', and ''global significance of 2.6$\sigma$ ... maximum local significance of 3.4$\sigma$''. The arXiv API itself returns 'not found' for this ID, so the quotes come from abstract text only and are not cite()-validated.

\textbf{Finding:} Abstract-level literature anchors from OpenAlex. (a) arXiv:2403.19598 (S. P. Martin, PRD 109, 095045, 2024) says a thermal higgsino-like DM ''mass must be between about 100 GeV and 1150 GeV''. (b) arXiv:2607.20588 (Bisal et al., 2026) says ''NLO corrections lower the DM-nucleon cross section below the neutrino floor'' for pure higgsinos, and that a pure higgsino ''typically evades current direct detection limits''. This abstract quotes no numerical cross section, so the target of about 1e-48 to 1e-49 cm\textasciicircum{}2 remains unsourced. (c) arXiv:2410.17137 (XLZD Design Book, EPJC 2025) reaches the neutrino fog, with 3 sigma evidence potential down to 3e-49 cm\textasciicircum{}2 at 40 GeV. (d) Planck 2018 papers 1807.06205 (overview) and 1907.12875 (likelihoods) were found, but the abstract of 1807.06209 (parameters) was not retrieved, so Omega\_c h\textasciicircum{}2 = 0.120 +/- 0.001 is still unquoted.

\textbf{Blocked --- Verify and cite the key literature: the higgsino relic crossing mass, the LZ 2.84 t-yr extended-window result, the solar-capture delta bound, the elastic loop-level cross section relative to the XENONnT limits and the neutrino fog, and the Planck Omega h\textasciicircum{}2. Obtain verbatim quotes..} (The arXiv API behind get\_arxiv\_metadata, both download tools and cite's metadata check returns 'not found', HTTP 429 or timeouts even for 1807.06209, so no full text could be read and no quote could be validated with cite(). Only OpenAlex abstracts are reachable.) No cite() entries could be made, so the success criterion of at least 5 cited quotes was not met. The arXiv-API layer is still down. get\_arxiv\_metadata returned 'not found' for 2609.02823, 2609.02775, 1807.06209 and 2410.15733, including the well-established Planck ID, so these results say nothing about whether the papers exist. download\_arxiv\_full\_text failed for 1807.06209, download\_full\_arxiv\_paper crashed with '''Search' object has no attribute 'results''', and cite() failed its metadata check for 2410.15733 (tc-...-64, -73, -77, -80, -81, -82, -84; dead end x-02). Only the OpenAlex-backed search\_arxiv returned anything, and it gives abstracts only.

Reference table from OpenAlex abstracts. None of these quotes are validated by cite().

\begin{itemize}
  \item arXiv:2609.02823 (LZ, 2026-09-02): the OpenAlex record exists and its abstract text matches the board. Exposure 2.84 tonne-years, recoil window up to about 270 keV, one event at 248 +/- 23 (stat) +/- 23 (sys) keV, 2.6 sigma global and 3.4 sigma local significance (f-013, tc-...-85). The elastic SI limit at 1 TeV is not in the abstract.
  \item arXiv:2403.19598 (Martin, PRD 109, 095045): a thermal higgsino-like DM mass must be between about 100 and 1150 GeV (f-014, tc-...-86).
  \item arXiv:2410.15733: the abstract quote 'Higgsino mass of \textasciitilde{}1.2 TeV' was already recorded by d-04 (f-007). cite() failed on it.
  \item arXiv:2607.20588 (2026): next-to-leading-order corrections push the pure-higgsino nucleon cross section below the neutrino floor. The abstract gives no number, so the value of about 1e-48 to 1e-49 cm\textasciicircum{}2 is still unsourced.
  \item arXiv:2410.17137 (XLZD): neutrino-fog reach, with 3 sigma evidence potential at 3e-49 cm\textasciicircum{}2 for a 40 GeV WIMP.
  \item Planck 2018: search\_arxiv found 1807.06205 and 1907.12875, but 1807.06209 did not appear in the results and its metadata lookup returned 'not found'. The value Omega\_c h\textasciicircum{}2 = 0.120 +/- 0.001 remains unquoted (tc-...-83).
\end{itemize}
Not checked, for lack of tool budget: 2609.02775 (its metadata lookup returned 'not found'), 2609.02608, 2609.26698, 2609.15321 and 2610.00226. The solar-capture bound delta \textgreater{} 566 keV is therefore still unquoted.

\textbf{Critic, wave 2} --- continue: 2 findings stand up to the evidence, 4 questioned ('The indirect-detection status of the about 1.1 TeV thermal higgsino, taken from abstracts ': The quoted excerpts do not contain the claimed values. tc-63 is cut off at the 2608.07234 author list, so the 546 h exposure, the 2.3e-28 cm\textasciicircum{}3/s Einasto line limit at 1 TeV and the 'probed for DM Milky Way models' quote are not visible. tc-71 shows only 2609.01590. Neither 2405.13104 ('within an order one factor', H.E.S.S. subleading) nor 2210.03140 ('out of reach by at least a factor of a few') appears in any cited excerpt. The 'not excluded' conclusion is plausible, but it needs to be re-posted with the abstract text that actually holds these values.; 'CTAO is the decisive future test. arXiv:2008.00692 (Rinchiuso et al., PRD 103, 023011) fin': The excerpts show only the title and authors of 2008.00692 and the start of the 2609.01590 entry. The quoted CTA sentence, the 2405.13104 'at high significance... for all but the most pessimistic DM profiles' quote and the 2506.08084 CTAO-North projection are not in the evidence. The physics point that the server's p-wave sigma\_v of about 1e-41 cm\textasciicircum{}3/s is wrong for an s-wave, Sommerfeld-enhanced higgsino is sound. Even so, confidence 'high' is not justified for abstract-level quotes that cannot be seen in the excerpts.; 'The OpenAlex record for arXiv:2609.02823 (LZ, 'Search for dark matter particle interaction': tc-85 is cut off after the LZ author list. The excerpt contains the ID and title but not the date, the DOI, '2.84 tonne-years', '270 keV', '248+/-23 (stat)+/-23 (sys) keV' or the 2.6/3.4 sigma significances. tc-64 and tc-81 only document API failures. The quoted abstract sentences need to be in the excerpt itself.; 'Abstract-level literature anchors from OpenAlex. (a) arXiv:2403.19598 (S. P. Martin, PRD 1': Only part (d) is supported: tc-77 and tc-82 show 1807.06209 'not found', and tc-83 shows 1907.12875. Parts (a) to (c) are not. The visible part of tc-86 shows 2203.02309, not 2403.19598, 2607.20588 or 2410.17137. The Martin 'between about 100 GeV and 1150 GeV' quote, the NLO neutrino-floor statement and the XLZD 3e-49 cm\textasciicircum{}2 figure are not in any quoted excerpt.).

Issues: The (794/776)\textasciicircum{}2 rescaling of delta(N=1) and delta(N=0.3) from 352.5/356.4 keV to about 369/373 keV is a heuristic. It assumes the whole N(delta) curve stretches with v\_max\textasciicircum{}2, but near the kinematic edge the count depends on how the halo tail is populated between v\_min and v\_max, not only on where the edge sits. The arithmetic checks out (mu = 109.8 GeV gives delta\_max = 367.8 and 385 keV), but the shifted window has not been computed. The claim that 2609.02608 uses v\_max = 794 km/s is also unsourced, so attributing 17 keV of the 24.5 keV offset to the cutoff is an assumption.; The 350 keV point (N = 1.97, peak E\_R = 315 keV) and the 300 keV point (744.6) come from wave-1 calls under d-02 and are not among this finding's evidence ids. The interpolations check out: ln-linear between 1.97 and 0.494 gives 352.45 keV, and between 0.494 and 0.085 gives 356.4 keV. The finding should still cite the d-02 call ids for the 350 keV anchor.; The board brief still says 'I did not compute any relic densities... No inelastic-DM or relic-density tools were available'. This contradicts d-01, which computed Omega h\textasciicircum{}2 at five masses (crossing at about 1158 GeV), and d-02 and d-05, which computed LZ event counts. The brief must be rewritten before synthesis. It should also downgrade the 'no Fermi-LAT exclusion' parts of the f-005 viable count, given the conclusion that f-002 is superseded..

Stopping after wave 2: budget reached (2/2 waves used). Writing the report from the board.

\textit{Research report:}\par
\subsubsection*{The 1.0--1.2 TeV thermal higgsino against LZ, XENONnT, Planck and gamma-ray limits: toy-model calculations and an abstract-level literature survey}
\subsubsection*{Abstract}
This investigation asked where the pure higgsino thermal relic, at about 1.0--1.2 TeV, sits relative to current limits.

\textbf{Tool calculations.} An inelastic-dark-matter (inelastic-DM) server was used for relic and recoil calculations.

\begin{itemize}
  \item \textbf{Relic density.} The server's relic model is a single-channel freeze-out scaling calibrated near 1.1 TeV. It reaches the Planck abundance at about 1158 GeV. This lies inside the 1.08--1.2 TeV literature range, but it is a calibration result, not a Sommerfeld/coannihilation prediction.
  \item \textbf{LZ inelastic events.} The server's Z-exchange higgsino predicts these LZ counts at 1.1 TeV for the 2.84 tonne-year extended nuclear-recoil window:
\end{itemize}
\begin{itemize}
  \item \textbf{Kinematic edge and one-event window.} The cutoff sits at about 368 keV for the tool's 776 km/s halo cutoff. It would sit at about 385 keV for 794 km/s. The ''about one event'' window is only about 4 keV wide, at 352--356 keV.
\end{itemize}
\begin{longtable}{@{}p{0.480\textwidth}p{0.480\textwidth}@{}}
\toprule
\textbf{Mass splitting delta (keV)} & \textbf{Expected LZ events} \\ \midrule \endhead
300 & 744.6 (excluded) \\
350 & about 2 \\
380 and above & 0 \\
\bottomrule\end{longtable}
\textbf{Literature, from abstracts only.} None of the following could be verified against full texts.

\begin{itemize}
  \item Solar-capture limits from IceCube require delta \textgreater{} 566 keV.
  \item That puts every solar-capture-allowed thermal higgsino beyond the inelastic reach of xenon.
  \item The pure-higgsino elastic cross section is described as lying below the neutrino floor.
  \item The higgsino is not excluded by current gamma-ray line searches, and the Cherenkov Telescope Array Observatory (CTAO) is identified as the decisive future test.
\end{itemize}
\textbf{Bottom line.} The full-abundance thermal higgsino at about 1.1 TeV is not excluded. Its interpretation as the source of the sub-MeV-splitting LZ 248 keV event is disfavoured. With splittings of order an MeV it is effectively invisible to current xenon detectors.

\subsubsection*{Background}
\textbf{Thermal mass.} With Sommerfeld enhancement and coannihilation included, a pure higgsino gives the full Planck abundance at roughly 1.0--1.2 TeV. Quoted values are:

\begin{itemize}
  \item about 1.08--1.1 TeV (arXiv:2609.02775, arXiv:2609.02608);
  \item ''\textasciitilde{}1.2 TeV'' (arXiv:2410.15733);
  \item an upper end of about 1150 GeV, with the thermal higgsino-like mass required to lie ''between about 100 GeV and 1150 GeV'' (arXiv:2403.19598).
\end{itemize}
Lighter thermal higgsinos are under-abundant by a factor of 5--20 in natural SUSY (arXiv:1803.11210).

\textbf{Direct detection.} The pure higgsino is inelastic dark matter. Z-mediated scattering requires the neutral-state splitting delta to be no more than a few hundred keV. If delta lies beyond the kinematic reach of xenon, only loop-level elastic scattering remains, and that rate sits near or below the neutrino fog (arXiv:2609.01504, arXiv:2410.02723). One 2026 abstract states that next-to-leading-order corrections push the pure-higgsino nucleon cross section ''below the neutrino floor'' (arXiv:2607.20588). Higgsino--gaugino mixing raises the elastic rate (arXiv:2606.31013).

\textbf{The 2026 LZ result.} LZ reported a 2.84 tonne-year search with a recoil window extended to about 270 keV (arXiv:2609.02823). It found one event at 248 $\pm$ 23 (stat) $\pm$ 23 (sys) keV, with 2.6$\sigma$ global and 3.4$\sigma$ local significance.

Interpretations of that event as a higgsino give these preferred splittings:

\begin{longtable}{@{}p{0.320\textwidth}p{0.320\textwidth}p{0.320\textwidth}@{}}
\toprule
\textbf{Assumption or method} & \textbf{Preferred delta} & \textbf{Source} \\ \midrule \endhead
Standard Halo Model & about 350 keV & arXiv:2609.01504 \\
Standard Halo Model & 377 keV & arXiv:2609.02608 \\
Sideband optimum at 1 TeV & 339 keV, about 1.9$\sigma$ tension & arXiv:2609.26698 \\
Including the Large Magellanic Cloud (LMC) & about 480--520 keV & arXiv:2609.01504; 481--519 keV in arXiv:2609.01590 \\
\bottomrule\end{longtable}
\textbf{Tensions with the event interpretation.}

\begin{itemize}
  \item \textbf{Solar capture.} IceCube non-observation requires delta \textgreater{} 566 keV (arXiv:2609.02775). This is supported by arXiv:2609.11833 and arXiv:2609.07807. arXiv:2609.06760 quotes 0.51--0.56 MeV.
  \item \textbf{Empty sideband.} The absence of events above 350 keV is in tension with the thermal-mass fit (arXiv:2609.09385, arXiv:2609.26698).
  \item \textbf{Halo uncertainties.} These cannot close the gap (arXiv:2609.15321). Under the Standard Halo Model, delta \textgreater{} 566 keV would need v0 \textgreater{} 240 km/s. With the LMC included, the largest workable splitting is 514 keV.
\end{itemize}
\textbf{Surviving branch.} A Planck-compatible thermal branch survives at delta of about 0.85--1.34 MeV (arXiv:2610.00226). arXiv:2609.09136 proposes a non-higgsino alternative.

\textbf{Caveat on sources.} All 2026-dated identifiers came from an OpenAlex-backed search and could not be resolved through the arXiv API. They should be checked on arxiv.org before anyone cites them.

\subsubsection*{Results}
\paragraph{Relic density from the inelastic-DM server (toy model)}
The relic tool was run with the higgsino flag set. This flag sets $\sigma$v = 2.5$\times$10$^-$$^2$? cm$^3$/s $\times$ (1100 GeV/m)$^2$ and feeds it into a Kolb--Turner freeze-out estimate, with x\_f about 24.5 and g* = 86.25. The tool has no Sommerfeld or coannihilation switch.

\begin{itemize}
  \item \textbf{Planck crossing.} Linear interpolation reaches $\Omega$ h$^2$ = 0.120 at about 1158 GeV. The server's $\pm$0.0012 band gives $\pm$6 GeV; the $\pm$0.001 band gives about $\pm$5 GeV. A direct call at 1159 GeV returns 0.120.
  \item \textbf{Comparison with the literature.} This crossing is inside the 1.08--1.2 TeV range. It is about 60--80 GeV above the 1.08--1.1 TeV values and close to the \textasciitilde{}1.2 TeV value of arXiv:2410.15733.
  \item \textbf{Interpretation.} The agreement reflects how the tool's cross section is normalized. It should be read as a toy-model crossing, not a higgsino prediction.
  \item \textbf{Tests not run.} The with/without-Sommerfeld comparison and splitting-dependence tests could not be done.
  \item \textbf{Pass flag.} The tool marks results as ''consistent'' with Planck even at 7--25$\sigma$ tension; for example, 1100 GeV carries 9.8$\sigma$ tension. This looks like a bug.
\end{itemize}
\paragraph{LZ inelastic event counts and the kinematic reach of xenon}
\begin{longtable}{@{}p{0.480\textwidth}p{0.480\textwidth}@{}}
\toprule
\textbf{m\_$\chi$ (GeV)} & \textbf{$\Omega$ h$^2$} \\ \midrule \endhead
900 & 0.0731 \\
1000 & 0.0898 \\
1100 & 0.1083 \\
1200 & 0.1285 \\
1300 & 0.1503 \\
\bottomrule\end{longtable}
The recoil tool uses:

\begin{itemize}
  \item Z exchange, with $\sigma$\_n = 7.43$\times$10$^-$$^3$? cm$^2$ per nucleon;
  \item the LZ 2.84 tonne-year extended nuclear-recoil window;
  \item a hard halo cutoff at v\_esc + v\_E = 776 km/s.
\end{itemize}
Expected LZ counts at m\_$\chi$ = 1100 GeV:

\begin{longtable}{@{}p{0.240\textwidth}p{0.240\textwidth}p{0.240\textwidth}p{0.240\textwidth}@{}}
\toprule
\textbf{delta (keV)} & \textbf{Expected LZ events} & \textbf{Spectrum peak (keV)} & \textbf{Notes} \\ \midrule \endhead
300 & 744.6 & not recorded & Grossly excluded against 1 observed event \\
350 & 1.97 & 315 & P(N $\geq$ 1) = 0.86; the 248 keV event needs v\_min = 762 km/s, which is accessible \\
380 & 0.00 & 342 & The 248 keV event needs v\_min = 799 km/s, which is inaccessible \\
500 & 0.00 & 450 &  \\
566 & 0.00 & 510 &  \\
700 & 0.00 & 630 &  \\
\bottomrule\end{longtable}
\textbf{Hand check of the edge.} The maximum splitting is $\delta$\_max = $\mu$ v\_max$^2$/2. With m\_Xe of about 122 GeV, the reduced mass is $\mu$ = 109.8 GeV.

\begin{itemize}
  \item For v\_max = 776 km/s, $\delta$\_max is about 368 keV.
  \item For v\_max = 794 km/s (544 + 250 km/s), $\delta$\_max is about 385 keV.
\end{itemize}
This matches the tool, which gives events at 350 keV and none at 380 keV.

\textbf{Mass dependence.} The edge depends only weakly on mass. At delta = 380 keV, v\_min for the 248 keV event is 802, 799 and 796 km/s at 1000, 1100 and 1200 GeV, all above 776 km/s.

\textbf{Consequence.} The solar-capture bound of delta \textgreater{} 566 keV from the literature therefore puts every solar-capture-allowed thermal higgsino beyond LZ's inelastic reach under the Standard Halo Model.

\textbf{Annual modulation.} At 1100 GeV, 350 keV and E\_R = 248 keV, the tool gives:

\begin{itemize}
  \item a fractional amplitude of 98.5\%, peaking on June 2;
  \item 22.5\% of the peak rate on the event date of September 1.
\end{itemize}
A recoil-spectrum plot for 1100 GeV and 380 keV was written to \texttt{/\allowbreak{}Users/\allowbreak{}nramachandra/\allowbreak{}github-integration/\allowbreak{}inelastic-dm-server/\allowbreak{}output/\allowbreak{}recoil\_\allowbreak{}spectrum.png}. No numbered figure was produced for this report.

\paragraph{The one-event splitting window}
A finer scan at 1100 GeV gives:

\begin{itemize}
  \item \textbf{Window location.} Log-linear interpolation gives delta(N=1) = 352.5 keV and delta(N=0.3) = 356.4 keV. The window is therefore about 4 keV wide, and N falls by about 4--6 times per 5 keV.
  \item \textbf{Where the event becomes inaccessible.} The 248 keV event drops out between 360 and 365 keV, where v\_min rises from 774 to 781 km/s.
  \item \textbf{Thermal mass.} At the tool's thermal mass of 1159 GeV, delta = 360 keV gives N = 0.135, against 0.085 at 1100 GeV. This moves delta(N=1) up by only about 1--2 keV.
  \item \textbf{Cause of the zero counts.} The zero counts at delta $\geq$ 370 keV come from the hard 776 km/s cutoff, not the energy window.
  \item \textbf{Halo dependence (not computed).} The recoil tool has no halo inputs, so halo dependence could not be computed. An analytic rescaling by (794/776)$^2$ = 1.047 would move the window to about 369--373 keV. This is a heuristic: near the edge, the count depends on how the velocity tail is populated, not only on where the edge sits.
\end{itemize}
\begin{longtable}{@{}p{0.320\textwidth}p{0.320\textwidth}p{0.320\textwidth}@{}}
\toprule
\textbf{delta (keV)} & \textbf{Expected LZ events} & \textbf{Spectrum peak (keV)} \\ \midrule \endhead
350 & 1.97 & 315 \\
355 & 0.494 & 320 \\
360 & 0.085 & 324 \\
365 & 0.009 & 329 \\
370 & 0.00 & 333 \\
375 & 0.00 & 338 \\
\bottomrule\end{longtable}
Comparison with literature fits:

\begin{itemize}
  \item \textbf{350 keV (arXiv:2609.01504).} Agrees with the tool directly.
  \item \textbf{377 keV (arXiv:2609.02608).} The offset is 24.5 keV. About 17 keV of it would be explained if that paper used a 794 km/s cutoff, but this is an assumption, not a sourced fact. The remaining roughly 5--8 keV is unexplained. Possible causes are the halo tail shape, v\_E at the event date, and cross-section or efficiency normalization; none was tested.
  \item \textbf{339 keV (arXiv:2609.26698).} This is a normalization-free sideband-ratio optimum. With the tool's normalization, splittings this low give hundreds of events, so it is not a like-for-like comparison.
\end{itemize}
The tool therefore does not by itself account for the spread among literature fits.

\paragraph{Parameter scan and constraint checker}
A 17 $\times$ 17 scan covered m\_$\chi$ = 800--1400 GeV and delta = 0--800 keV. It reported 144 of 289 points ''viable'' and 9 points predicting 0.3--5 LZ events. The best-fit point is m\_$\chi$ = 988 GeV, delta = 350 keV. Output files are:

\begin{itemize}
  \item \texttt{/\allowbreak{}Users/\allowbreak{}nramachandra/\allowbreak{}github-integration/\allowbreak{}inelastic-dm-server/\allowbreak{}output/\allowbreak{}scan\_\allowbreak{}idm.npz}
  \item \texttt{/\allowbreak{}Users/\allowbreak{}nramachandra/\allowbreak{}github-integration/\allowbreak{}inelastic-dm-server/\allowbreak{}output/\allowbreak{}constraint\_\allowbreak{}summary.png}
\end{itemize}
These numbers should not be taken as physics, for three reasons:

\begin{itemize}
  \item The viable count uses the relic pass flag that stays true at 7--25$\sigma$.
  \item It uses an indirect-detection model that is wrong for a higgsino (see below).
  \item The best fit fits the LZ event, not the relic. At 988 GeV the tool's own relic curve gives $\Omega$ h$^2$ of about 0.087, so this point is under-abundant.
\end{itemize}
The constraint checker rated four points VIABLE: 1100 GeV with delta = 400, 566 and 1000 keV, and 1159 GeV with 400 keV. All four have N\_LZ = 0.00 and no Fermi-LAT or CMB exclusion. Its present-day $\sigma$v of about 7$\times$10$^-$?$^2$ to 5$\times$10$^-$?$^1$ cm$^3$/s is labelled ''p-wave + mass-splitting suppressed''. That is inappropriate for an s-wave, Sommerfeld-enhanced higgsino, whose cross section is about 10$^-$$^2$? cm$^3$/s.

\paragraph{Elastic scattering (literature only)}
The XENONnT 3.1 tonne-year search (arXiv:2502.18005) reports:

\begin{itemize}
  \item a minimum spin-independent limit of 1.7$\times$10$^-$?? cm$^2$ at 30 GeV (90\% CL);
  \item a best median sensitivity of 1.4$\times$10$^-$?? cm$^2$ at 41 GeV.
\end{itemize}
The earlier 1.09 tonne-year result (arXiv:2303.14729) reached 2.58$\times$10$^-$?? cm$^2$ at 28 GeV.

Neither abstract gives the limit at 1 TeV. The figure of about 10$^-$?? cm$^2$ is only an extrapolation using linear high-mass scaling.

XLZD (arXiv:2410.17137) reaches the neutrino fog, with 3$\sigma$ evidence potential down to 3$\times$10$^-$?? cm$^2$ at 40 GeV.

The pure-higgsino loop-level elastic cross section of about 10$^-$??--10$^-$?? cm$^2$ was not sourced. arXiv:2607.20588 places it qualitatively below the neutrino floor without giving a number.

\paragraph{Indirect detection (literature only)}
These findings are based on abstracts only. The full text of arXiv:2405.13104 could not be downloaded.

\begin{itemize}
  \item \textbf{Fermi-LAT.} A 15-year line-like search comes ''within an order one factor'' of the sensitivity needed to detect the higgsino, for some simulated Milky Way profiles. In the same analysis, H.E.S.S. is subleading to Fermi (arXiv:2405.13104).
  \item \textbf{Current imaging Cherenkov telescopes.} At their ultimate sensitivity, H.E.S.S.-like instruments leave the higgsino ''out of reach by at least a factor of a few''. The dominant systematic is the dark-matter density in the inner Galaxy (arXiv:2210.03140).
  \item \textbf{H.E.S.S. Inner Galaxy Survey.} This 2026-dated, unverified paper used 546 h of data (arXiv:2608.07234). It quotes a line limit of 2.3$\times$10$^-$$^2$? cm$^3$/s at 1 TeV for an Einasto profile and says the thermal higgsino is ''probed'' for the first time. Its abstract does not say whether, or for which profiles, the higgsino is excluded.
  \item \textbf{MAGIC.} No MAGIC-specific limit surfaced.
  \item \textbf{Future instruments.}
  \item CTA could be the first experiment with sensitivity to the thermal higgsino (arXiv:2008.00692).
  \item CTA and SWGO could detect it ''at high significance'' for all but the most pessimistic profiles (arXiv:2405.13104).
  \item CTAO-North could reach the required sensitivity within a decade (arXiv:2506.08084).
\end{itemize}
The predicted line cross section at 1.1 TeV, and its ratio to current limits, could not be extracted.

\paragraph{Literature verification status}
Every arXiv-API-backed tool failed: metadata lookup, both download tools, and the metadata check inside the citation tool. They returned ''not found'', HTTP 429 or timeouts, even for the Planck 2018 parameters paper (1807.06209). No quote could be validated and no bibliography entry was made.

Abstract-level text was obtained through OpenAlex for:

\begin{itemize}
  \item the LZ result (arXiv:2609.02823);
  \item arXiv:2403.19598;
  \item arXiv:2410.15733;
  \item arXiv:2607.20588;
  \item arXiv:2410.17137.
\end{itemize}
Planck overview and likelihood papers (1807.06205, 1907.12875) were found. The parameters abstract was not, so $\Omega$\_c h$^2$ = 0.120 $\pm$ 0.001 is used here as the standard value but remains unquoted.

\subsubsection*{Discussion}
Answering ''where does the higgsino sit?'' requires assembling three pieces.

\textbf{First, the relic target.} This is set by the literature at about 1.1 TeV, give or take about 0.1 TeV. The tool's crossing at about 1158 GeV is consistent with that only because the tool is calibrated there.

\textbf{Second, inelastic direct detection.} The xenon kinematics are robust and do not depend on the toy model. At the thermal mass, $\delta$\_max is about 368--385 keV depending on the velocity cutoff, nearly independent of mass. The sub-MeV LZ interpretation therefore requires a splitting in a narrow band just below this edge. In the tool, that band is 352--356 keV for one to 0.3 events, and lower splittings are grossly excluded.

The solar-capture bound reported in the literature (delta \textgreater{} 566 keV, or 0.51--0.56 MeV) lies far above this edge under the Standard Halo Model. It also lies above the roughly 480--520 keV LMC-tail values. This is why the LZ interpretation is disfavoured while the higgsino itself survives at delta of order an MeV (arXiv:2610.00226). In that regime, inelastic scattering is closed in xenon. The elastic rate is described as below the neutrino floor, although no number was sourced here.

\textbf{Third, indirect detection.} The literature consensus from abstracts is that the higgsino is within an order-one factor of current Fermi-LAT line sensitivity and is not excluded, with strong dependence on the halo profile. CTAO is the decisive future test.

The results agree with the literature on:

\begin{itemize}
  \item the location of the xenon edge;
  \item the roughly 350 keV preferred splitting under a Standard-Halo-like cutoff.
\end{itemize}
They are in tension, or at least unreconciled, with:

\begin{itemize}
  \item the 377 keV fit of arXiv:2609.02608, where 5--8 keV is unexplained even after the assumed cutoff correction;
  \item the 339 keV sideband optimum, which is not a comparable quantity.
\end{itemize}
Dead ends:

\begin{itemize}
  \item \textbf{Verbatim citations.} Verbatim citation of Planck 2018 VI, arXiv:2410.15733, arXiv:2609.02823 and arXiv:2609.02775 was tried twice. Each time, the arXiv API layer returned ''not found'', 429 errors or internal errors, even for long-established papers. A ''not found'' result therefore says nothing about whether a paper exists.
  \item \textbf{Tool-based Sommerfeld, halo and indirect-detection studies.} These were impossible because the server has no Sommerfeld switch and no halo parameters, and its indirect-detection model is p-wave.
\end{itemize}
\subsubsection*{Limitations}
\begin{itemize}
  \item \textbf{Relic model.} The relic calculation is a toy single-channel scaling. No Sommerfeld or coannihilation physics was computed, and no tool for it was available.
  \item \textbf{Halo model.} Recoil predictions use a single hard 776 km/s cutoff. Halo sensitivity, including v0, v\_esc, a smooth tail and the LMC, was only estimated by an analytic rescaling. The Standard Halo Model v0 scan of arXiv:2609.15321 and the sideband analysis of arXiv:2609.26698 were not reproduced.
  \item \textbf{Binned counts.} Counts in the [14,225), [225,271] and (271,800] keV bins were not obtained. The tool returns integrated rates only.
  \item \textbf{Constraint checker.} Its relic pass flag and its indirect-detection model are unreliable. Its viability counts and ''no Fermi-LAT exclusion'' verdicts should be disregarded.
  \item **
\end{itemize}
\end{document}
